% ============================================================================
% SEC05.TEX - Section 11.5: Line of Sight Detection
% ============================================================================

\section{Line of Sight Detection}

\begin{learningoutcomes}
\item Mengimplementasikan line of sight detection
\item Menggunakan Raycast untuk deteksi
\item Membuat field of view untuk AI
\end{learningoutcomes}

\subsection{Line of Sight dengan Raycast}

\begin{lstlisting}[language={[Sharp]C}]
public class LineOfSight : MonoBehaviour
{
    public Transform target;
    public float viewDistance = 10f;
    public float viewAngle = 45f;
    public LayerMask obstacleLayer;
    
    public bool CanSeeTarget()
    {
        Vector3 direction = (target.position - transform.position).normalized;
        float angle = Vector3.Angle(transform.forward, direction);
        
        if (angle > viewAngle / 2f)
            return false;
        
        float distance = Vector3.Distance(transform.position, target.position);
        if (distance > viewDistance)
            return false;
        
        RaycastHit hit;
        if (Physics.Raycast(transform.position, direction, out hit, viewDistance, obstacleLayer))
        {
            if (hit.collider.transform == target)
                return true;
        }
        
        return false;
    }
}
\end{lstlisting}

\begin{catatan}
Line of sight membuat AI lebih realistis - hanya mendeteksi player jika dalam field of view dan tidak terhalang obstacle.
\end{catatan}

\subsection{Ringkasan Section}

\begin{itemize}
\item Line of sight menggunakan Raycast dan angle check
\item Field of view menentukan sudut deteksi
\item Obstacle layer untuk menghindari deteksi melalui dinding
\item Membuat AI lebih realistis dan fair
\end{itemize}
